\documentclass{article}
\usepackage{amsmath,amssymb,amsthm}
\usepackage{algorithm}
\usepackage[noend]{algorithmic}
\usepackage{enumitem}
\usepackage{hyperref}
\newtheorem{theorem}{Theorem}
\newtheorem{lemma}{Lemma}
\newtheorem{assumption}{Assumption}
\newtheorem{remark}{Remark}
\newcommand{\prox}{\operatorname{prox}}
\newcommand{\R}{\mathbb{R}}
\newcommand{\E}{\mathbb{E}}

\begin{document}

\section*{Proximal Gradient Method (Forward–Backward Splitting)}

\section*{Mathematical Formulation}
We consider the composite convex optimization problem  

\[
\min_{x\in\R^{d}}\; \Phi(x)\qquad\text{with}\qquad 
\Phi(x):=f(x)+g(x),
\]

where  

\begin{itemize}[nosep]
\item $f:\R^{d}\to\R$ is convex, continuously differentiable and its gradient is $L$–Lipschitz, i.e.  
\[
\|\nabla f(x)-\nabla f(y)\|\le L\|x-y\|\quad\forall x,y\in\R^{d};
\]
\item $g:\R^{d}\to\R\cup\{+\infty\}$ is convex, proper and lower‑semicontinuous (possibly non‑smooth);
\item $\eta>0$ is a stepsize (also called learning rate) satisfying $0<\eta\le 1/L$.
\end{itemize}

The proximal operator of $g$ with parameter $\eta$ is  

\[
\prox_{\eta g}(z):=\arg\min_{u\in\R^{d}}\Bigl\{ g(u)+\frac{1}{2\eta}\|u-z\|^{2}\Bigr\}.
\]

The algorithm generates a sequence $\{x_{k}\}_{k\ge0}$ by  

\[
\boxed{\;x_{k+1}= \prox_{\eta g}\bigl(x_{k}-\eta\nabla f(x_{k})\bigr)\;},\qquad k=0,1,2,\dots
\tag{1}
\]

\section*{Pseudocode}
\begin{algorithm}[H]
\caption{Proximal Gradient Method}
\begin{algorithmic}[1]
\REQUIRE Initial point $x_{0}\in\R^{d}$, stepsize $\eta\in(0,1/L]$.
\FOR{$k=0,1,2,\dots$}
    \STATE Compute the gradient $g_{k}:=\nabla f(x_{k})$.
    \STATE Perform a forward step: $y_{k}=x_{k}-\eta g_{k}$.
    \STATE Perform a backward (proximal) step: $x_{k+1}= \prox_{\eta g}(y_{k})$.
    \IF{termination criterion satisfied (e.g., $\|x_{k+1}-x_{k}\|\le \varepsilon$)}
        \STATE \textbf{break}
    \ENDIF
\ENDFOR
\RETURN $x_{k+1}$
\end{algorithmic}
\end{algorithm}

\section*{Dry Run Example}
We minimise the scalar function  

\[
f(w)=(w-3)^{2},\qquad g\equiv0,
\]
so that $\Phi(w)=f(w)$.  
The gradient is $\nabla f(w)=2(w-3)$.  
Choose $x_{0}=0$ and stepsize $\eta=0.1$ (note that $L=2$, thus $\eta\le1/L=0.5$).

Since $g\equiv0$, $\prox_{\eta g}(z)=z$.  The iteration (1) reduces to  

\[
x_{k+1}=x_{k}-\eta\,2(x_{k}-3).
\]

\begin{center}
\begin{tabular}{c|c|c|c}
$k$ & $x_{k}$ & $\nabla f(x_{k})$ & $\Phi(x_{k})=(x_{k}-3)^{2}$\\\hline
0 & $0$ & $-6$ & $9$\\
1 & $0-0.1(-6)=0.6$ & $2(0.6-3)=-4.8$ & $(0.6-3)^{2}=5.76$\\
2 & $0.6-0.1(-4.8)=1.08$ & $2(1.08-3)=-3.84$ & $(1.08-3)^{2}=3.6864$\\
3 & $1.08-0.1(-3.84)=1.464$ & $2(1.464-3)=-3.072$ & $(1.464-3)^{2}=2.3690$
\end{tabular}
\end{center}
The objective value decreases monotonically, illustrating the descent property of the method.

\newpage
\section*{Assumptions}
\begin{assumption}[Convexity and Smoothness of $f$]\label{ass:f}
$f$ is convex and differentiable on $\R^{d}$ with $L$‑Lipschitz continuous gradient:
\[
\|\nabla f(x)-\nabla f(y)\|\le L\|x-y\|,\qquad\forall x,y\in\R^{d}.
\]
\end{assumption}
\begin{assumption}[Convexity of $g$]\label{ass:g}
$g$ is proper, convex and lower‑semicontinuous. Its proximal operator $\prox_{\eta g}$ is well defined for any $\eta>0$.
\end{assumption}
\begin{assumption}[Stepsize]\label{ass:step}
The stepsize $\eta$ satisfies $0<\eta\le 1/L$.
\end{assumption}

\section*{Main Convergence Theorem}
\begin{theorem}[Sublinear convergence of the proximal gradient method]\label{thm:main}
Let $\{x_{k}\}$ be generated by \eqref{1} under Assumptions \ref{ass:f}–\ref{ass:step} and denote by $x^{\star}\in\arg\min\Phi$ any optimal solution. Then for all $k\ge1$  

\[
\Phi(x_{k})-\Phi(x^{\star})\;\le\;\frac{\|x_{0}-x^{\star}\|^{2}}{2\eta\,k}.
\tag{2}
\]

Consequently $\Phi(x_{k})\to\Phi(x^{\star})$ and the sequence $\{x_{k}\}$ is bounded.  
If, in addition, $f$ is $\mu$‑strongly convex ($\mu>0$), the method enjoys a linear rate  

\[
\|x_{k}-x^{\star}\|\le\bigl(1-\eta\mu\bigr)^{k}\|x_{0}-x^{\star}\|.
\tag{3}
\]
\end{theorem}

\section*{Theoretical Properties}
\begin{itemize}[nosep]
\item \textbf{Non‑expansiveness of the proximal operator.} For any $u,v\in\R^{d}$,
\[
\|\prox_{\eta g}(u)-\prox_{\eta g}(v)\|\le\|u-v\|.
\tag{4}
\]
This property is the cornerstone of the proof and guarantees that the backward step does not increase the distance generated by the forward (gradient) step.

\item \textbf{Stability w.r.t.\ stepsize.} The bound (2) holds for any $\eta\in(0,1/L]$; larger $\eta$ (up to $1/L$) yields a smaller constant $1/(2\eta)$, i.e. faster descent.

\item \textbf{Comparison to gradient descent.} When $g\equiv0$, $\prox_{\eta g}$ is the identity and (2) reduces to the classical $O(1/k)$ rate of gradient descent for convex smooth functions.

\item \textbf{Extension to stochastic or inexact gradients.} The same proof skeleton can be adapted by adding unbiased variance bounds, leading to $O(1/\sqrt{T})$ rates similar to SGD (see Example 1).

\end{itemize}

\section*{Discussion}
The proximal gradient method is a first‑order method that can handle a non‑smooth regulariser $g$ (e.g., $\ell_{1}$ norm) while preserving the simple gradient step for $f$. The convergence proof relies only on convexity, Lipschitz continuity of $\nabla f$, and the firmly non‑expansive nature of $\prox_{\eta g}$. No additional smoothness of $g$ is required.

\newpage
\section*{Preliminaries (Definitions and Lemmas)}
\begin{lemma}[Descent Lemma]\label{lem:descent}
If $\nabla f$ is $L$‑Lipschitz then for any $x,y\in\R^{d}$,
\[
f(y)\le f(x)+\langle\nabla f(x),y-x\rangle+\frac{L}{2}\|y-x\|^{2}.
\tag{5}
\]
\end{lemma}
\begin{proof}
Standard integration of the Lipschitz condition; see e.g. \cite{Nesterov2004}.
\end{proof}

\begin{lemma}[Firm non‑expansiveness of the proximal operator]\label{lem:firm}
For any $\eta>0$ and any $u,v\in\R^{d}$,
\[
\|\prox_{\eta g}(u)-\prox_{\eta g}(v)\|^{2}
\le \langle u-v,\prox_{\eta g}(u)-\prox_{\eta g}(v)\rangle .
\tag{6}
\]
\end{lemma}
\begin{proof}
From the optimality condition of the proximal sub‑problem,
\[
\frac{1}{\eta}(z-\prox_{\eta g}(z))\in\partial g\bigl(\prox_{\eta g}(z)\bigr),
\]
and monotonicity of the sub‑differential of a convex function, we obtain (6).
\end{proof}

\section*{Main Proof (Step‑by‑Step Derivation)}
Let $x^{\star}\in\arg\min\Phi$.  Define the auxiliary point  

\[
y_{k}=x_{k}-\eta\nabla f(x_{k})\qquad\text{(forward step).}
\tag{7}
\]

The next iterate is $x_{k+1}= \prox_{\eta g}(y_{k})$ (backward step).  

\begin{enumerate}
\item[\textbf{[Step 1]}]  Start from the squared distance to the optimum:
\[
\|x_{k+1}-x^{\star}\|^{2}
 =\| \prox_{\eta g}(y_{k})- \prox_{\eta g}(y^{\star})\|^{2},
\]
where $y^{\star}=x^{\star}-\eta\nabla f(x^{\star})$ (note that $x^{\star}= \prox_{\eta g}(y^{\star})$ because $0\in \nabla f(x^{\star})+\partial g(x^{\star})$).

\item[\textbf{[Step 2]}]  Apply the firm non‑expansiveness (Lemma~\ref{lem:firm}) with $u=y_{k}$ and $v=y^{\star}$:
\[
\|x_{k+1}-x^{\star}\|^{2}
\le \langle y_{k}-y^{\star},\,x_{k+1}-x^{\star}\rangle .
\tag{8}
\]

\item[\textbf{[Step 3]}]  Substitute $y_{k}=x_{k}-\eta\nabla f(x_{k})$ and $y^{\star}=x^{\star}-\eta\nabla f(x^{\star})$:
\[
\langle y_{k}-y^{\star},\,x_{k+1}-x^{\star}\rangle
= \langle x_{k}-x^{\star},\,x_{k+1}-x^{\star}\rangle
-\eta\langle\nabla f(x_{k})-\nabla f(x^{\star}),\,x_{k+1}-x^{\star}\rangle .
\tag{9}
\]

\item[\textbf{[Step 4]}]  Expand the first inner product:
\[
\langle x_{k}-x^{\star},\,x_{k+1}-x^{\star}\rangle
= \tfrac12\bigl(\|x_{k}-x^{\star}\|^{2}+\|x_{k+1}-x^{\star}\|^{2}
-\|x_{k+1}-x_{k}\|^{2}\bigr).
\tag{10}
\]

\item[\textbf{[Step 5]}]  Insert (10) into (8)–(9) and rearrange terms to isolate $\|x_{k+1}-x^{\star}\|^{2}$ on the left:
\[
\|x_{k+1}-x^{\star}\|^{2}
\le \|x_{k}-x^{\star}\|^{2}
-\|x_{k+1}-x_{k}\|^{2}
-2\eta\langle\nabla f(x_{k})-\nabla f(x^{\star}),\,x_{k+1}-x^{\star}\rangle .
\tag{11}
\]

\item[\textbf{[Step 6]}]  Use the convexity of $f$:
\[
f(x_{k})-f(x^{\star})\le \langle\nabla f(x_{k}),\,x_{k}-x^{\star}\rangle ,
\qquad
f(x_{k+1})-f(x^{\star})\le \langle\nabla f(x_{k+1}),\,x_{k+1}-x^{\star}\rangle .
\tag{12}
\]

\item[\textbf{[Step 7]}]  Add the two inequalities in (12) and subtract $\langle\nabla f(x_{k}),\,x_{k+1}-x^{\star}\rangle$ to obtain
\[
\langle\nabla f(x_{k})-\nabla f(x^{\star}),\,x_{k+1}-x^{\star}\rangle
\ge f(x_{k})-f(x^{\star}) -\langle\nabla f(x_{k}),\,x_{k+1}-x_{k}\rangle .
\tag{13}
\]

\item[\textbf{[Step 8]}]  Plug (13) into (11):
\[
\|x_{k+1}-x^{\star}\|^{2}
\le \|x_{k}-x^{\star}\|^{2}
-\|x_{k+1}-x_{k}\|^{2}
-2\eta\bigl[f(x_{k})-f(x^{\star})\bigr]
+2\eta\langle\nabla f(x_{k}),\,x_{k+1}-x_{k}\rangle .
\tag{14}
\]

\item[\textbf{[Step 9]}]  Since $x_{k+1}= \prox_{\eta g}(y_{k})$ with $y_{k}=x_{k}-\eta\nabla f(x_{k})$, we have the optimality condition  
\[
\frac{1}{\eta}\bigl(y_{k}-x_{k+1}\bigr)\in\partial g(x_{k+1}),
\]
which implies (by convexity of $g$)
\[
g(x_{k+1})-g(x^{\star})\le \frac{1}{\eta}\langle y_{k}-x_{k+1},\,x_{k+1}-x^{\star}\rangle .
\tag{15}
\]

\item[\textbf{[Step 10]}]  Expand the right‑hand side of (15) using $y_{k}=x_{k}-\eta\nabla f(x_{k})$:
\[
\frac{1}{\eta}\langle y_{k}-x_{k+1},\,x_{k+1}-x^{\star}\rangle
= \frac{1}{\eta}\langle x_{k}-x_{k+1},\,x_{k+1}-x^{\star}\rangle
-\langle\nabla f(x_{k}),\,x_{k+1}-x^{\star}\rangle .
\tag{16}
\]

\item[\textbf{[Step 11]}]  Rewrite the first term of (16) as
\[
\frac{1}{\eta}\langle x_{k}-x_{k+1},\,x_{k+1}-x^{\star}\rangle
= \frac{1}{2\eta}\bigl(\|x_{k}-x^{\star}\|^{2}
-\|x_{k+1}-x^{\star}\|^{2}
-\|x_{k+1}-x_{k}\|^{2}\bigr).
\tag{17}
\]

\item[\textbf{[Step 12]}]  Substitute (17) into (16) and then (16) into (15):
\[
g(x_{k+1})-g(x^{\star})
\le
\frac{1}{2\eta}\bigl(\|x_{k}-x^{\star}\|^{2}
-\|x_{k+1}-x^{\star}\|^{2}
-\|x_{k+1}-x_{k}\|^{2}\bigr)
-\langle\nabla f(x_{k}),\,x_{k+1}-x^{\star}\rangle .
\tag{18}
\]

\item[\textbf{[Step 13]}]  Add inequality (18) to inequality (14).  The terms $-\langle\nabla f(x_{k}),x_{k+1}-x^{\star}\rangle$ cancel, yielding
\[
\Phi(x_{k+1})- \Phi(x^{\star})
\le
\frac{1}{2\eta}\bigl(\|x_{k}-x^{\star}\|^{2}
-\|x_{k+1}-x^{\star}\|^{2}\bigr)
-\frac{1}{2}\|x_{k+1}-x_{k}\|^{2}.
\tag{19}
\]

\item[\textbf{[Step 14]}]  Drop the non‑negative term $-\frac12\|x_{k+1}-x_{k}\|^{2}$ to obtain a simpler inequality:
\[
\Phi(x_{k+1})- \Phi(x^{\star})
\le
\frac{1}{2\eta}\bigl(\|x_{k}-x^{\star}\|^{2}
-\|x_{k+1}-x^{\star}\|^{2}\bigr).
\tag{20}
\]

\item[\textbf{[Step 15]}]  Rearrange (20) and sum from $k=0$ to $k=K-1$:
\[
\sum_{k=0}^{K-1}\bigl[\Phi(x_{k+1})- \Phi(x^{\star})\bigr]
\le
\frac{1}{2\eta}\bigl(\|x_{0}-x^{\star}\|^{2}
-\|x_{K}-x^{\star}\|^{2}\bigr)
\le\frac{\|x_{0}-x^{\star}\|^{2}}{2\eta}.
\tag{21}
\]

\item[\textbf{[Step 16]}]  Since the left‑hand side is a sum of $K$ non‑negative terms, the smallest term is bounded by the average:
\[
\Phi(x_{K})- \Phi(x^{\star})
\le
\frac{1}{K}\sum_{k=0}^{K-1}\bigl[\Phi(x_{k+1})- \Phi(x^{\star})\bigr]
\le\frac{\|x_{0}-x^{\star}\|^{2}}{2\eta K}.
\tag{22}
\]

Equation (22) is exactly the rate (2) claimed in Theorem~\ref{thm:main}.  This completes the proof for the convex case. \hfill $\square$

\medskip
\noindent\textbf{Proof of the linear rate (3).}  If $f$ is $\mu$‑strongly convex, then for all $x$,
\[
f(x)\ge f(x^{\star})+\langle\nabla f(x^{\star}),x-x^{\star}\rangle+\frac{\mu}{2}\|x-x^{\star}\|^{2}.
\]
Since $0\in\nabla f(x^{\star})+\partial g(x^{\star})$, the optimality condition gives $-\nabla f(x^{\star})\in\partial g(x^{\star})$.  Using the same steps as above but now applying the strong convexity inequality to the term $f(x_{k})-f(x^{\star})$ leads to  

\[
\|x_{k+1}-x^{\star}\|^{2}\le (1-\eta\mu)\|x_{k}-x^{\star}\|^{2},
\]
and iterating yields (3).  ∎

\section*{Rate Derivation (With Constants)}
From (22) we have the explicit constant
\[
C:=\frac{\|x_{0}-x^{\star}\|^{2}}{2\eta},
\qquad\text{so that}\qquad
\Phi(x_{k})-\Phi(x^{\star})\le \frac{C}{k}.
\]
If the stepsize is chosen as the maximal admissible value $\eta=1/L$, then $C=\frac{L}{2}\|x_{0}-x^{\star}\|^{2}$.  Thus the bound depends linearly on the Lipschitz constant $L$ and quadratically on the distance of the initialization from the optimum.

\section*{Special Cases}
\begin{enumerate}[nosep]
\item \textbf{Pure gradient descent ($g\equiv0$).}  The proximal operator reduces to the identity, and (22) becomes the classic $O(1/k)$ rate for smooth convex functions.

\item \textbf{Lasso regularisation.}  For $g(x)=\lambda\|x\|_{1}$ the proximal step is soft‑thresholding:
\[
\bigl[\prox_{\eta g}(z)\bigr]_{i}= \operatorname{sign}(z_{i})\max\{|z_{i}|-\eta\lambda,0\}.
\]
All the previous analysis remains valid, giving the same $O(1/k)$ guarantee for the composite objective.

\item \textbf{Indicator of a convex set $C$.}  If $g(x)=\iota_{C}(x)$, then $\prox_{\eta g}$ is the Euclidean projection onto $C$.  The algorithm becomes projected gradient descent, and the same convergence bound holds.

\item \textbf{Strongly convex $f$ + convex $g$.}  The linear rate (3) applies, showing that the presence of a non‑smooth convex $g$ does not destroy exponential convergence as long as the smooth part is strongly convex.

\end{enumerate}

\section*{References}
\begin{thebibliography}{9}
\bibitem{Nesterov2004}
Y.~Nesterov, \emph{Introductory Lectures on Convex Optimization: A Basic Course}, Springer, 2004.
\end{thebibliography}

\end{document}